\section{Bonfire of the Vanities}

Savonarola and Machiavelli in Florence.

\begin{quotex}
Christianity has not been tried and found wanting; it has been found difficult and not tried. \flright{\textsc{Gilbert Chesterton}}

\end{quotex}
Chesterton must have failed to consider \textbf{Girolamo Savonarola} when he penned that line. First of all, there is no merit in a religion that is not worth trying. Second, every path is difficult, from the yoga path of \textbf{Patanjali}, to the eightfold path of \textbf{Buddha}, even the \textbf{Old Religion}. In all of them, there is devotion, duty, action, and wisdom. Only an egalitarian thinks that everyone must succeed the same way. Isn't that why we honor our Saints, Heroes, and Sages because some excel and point the way?

I won't repeat the history of Savanorola's rise to be Capo of Florence; those unfamiliar with that piece of history can easily look it up. In brief, Savanorola passed up the ``easy life" in his father's business. The bourgeois life of family and property had no appeal to him; but tradition did, so he chose the path of a Dominican monk. Driven by an intransigent piety, he was unaffected by the lures of money, sex, and power, much like \textbf{Dostoyevsky}'s idiot Prince. He was blessed [plagued?] with apocalyptic visions [hallucinations?] which came from God. These were the source of his power and his torment.

Initially, he did not stand out, but over time he became known for his devoutness and he attracted a devoted following to his fiery and inspiring sermons. A series of events including the death of Lorenzo de Medici, the invasion from France and the rebellion of Pisa led to a power vacuum which he was called upon to fill. Although Florence was at the peak of its economic power and the center of the intellectual Renaissance, this return to a Medieval cast of mind still appealed to a large portion of the public. As a sort of priest-king, he outlawed usury, made sodomy a capital crime (causing many in the upper crust to flee), and engaged roving street gangs to enforce moral codes. He is most known for the \textbf{Bonfire of the Vanities}, where citizens were encouraged to rid themselves of secular, sinful and vain items.

Was this truly the difficult Christian life Chesterton described? If so, was it truly that much different from life in the Ancient City? Evola regarded Savonarola as the last attempt of the ancient Aryan-Roman spirit to reassert itself amidst the decadence of the humanism of the Renaissance.

Eventually, Savonarola got in trouble with the Pope for pointing out the latter's immoralities and was driven from power; he, along with two younger monks, were burned. A young man in Florence, \textbf{Niccolo Machiavelli} was an acute observer of the rise and fall of Savonarola. As long as the people believed in him, Savonarola could rule; but when their faith faltered, he fell. Unburdened by a belief in the supernatural, Machiavelli wrote his impressions in the Prince.

\begin{quotex}
if we desire to discuss this matter thoroughly, to inquire whether these innovators can rely on themselves or have to depend on others: that is to say, whether, to consummate their enterprise, have they to use prayers or can they use force? In the first instance they always succeed badly, and never compass anything; but when they can rely on themselves and use force, then they are rarely endangered. Hence it is that all armed prophets have conquered, and the unarmed ones have been destroyed. Besides the reasons mentioned, the nature of the people is variable, and whilst it is easy to persuade them, it is difficult to fix them in that persuasion. And thus it is necessary to take such measures that, when they believe no longer, it may be possible to make them believe by force.

If Moses, Cyrus, Theseus, and Romulus had been unarmed they could not have enforced their constitutions for long — as happened in our time to Fra Girolamo Savonarola, who was ruined with his new order of things immediately the multitude believed in him no longer, and he had no means of keeping steadfast those who believed or of making the unbelievers to believe.

\end{quotex}
In \textit{The Man of the Renaissance}, Ralph Roeder offers up his own opinions. He writes regarding the failure of this experiment in Florence:

\begin{quotex}
Christianity and life were incompatible. What had goodness to do with life or compassion with nature? The only true solution was the renunciation of the world and the real core of his creed was death.

\end{quotex}
Roeder has identified the issue in spite of himself. Although Evola would hardly use the words ``goodness" and ``compassion", isn't he in some way close to Savonarola? However, to oppose nature, one must be hard, not compassionate. They agree that the true solution lies in renunciation, but that should lead to a higher, heroic life, not to death.

Yet Roeder follows Machiavelli in believing they are returning to paganism, although it is closer to the type of neopaganism that Evola opposes. Roeder makes this claim:

\begin{quotex}
In the cramped past, in the poverty and ignorance of the Middle Ages, it had been possible perhaps for so unworldly a faith to flourish; but with the affluence and culture the world had outgrown the ascetic faith of its forefathers and rediscovered, in the Classic Revival, that of its ancestors, for whom the first and last law was the satisfaction of life. Life — imperfect, ruthless, lusty, lawless — was richly enough, and the only mastery of nature lay in its imitation.

\end{quotex}
Yet, there is an inconsistency in Roeder when he writes of Machiavelli's view of the Ancient Romans:

\begin{quotex}
The secret of their success was their virtue, that virile virtu — courage, energy, skill, resourcefulness, strength — which was the power of a man to function efficiently and fulfill his purpose. And since patriotism was his purpose, whatever impeded it — humanity or personal scruples — must be sacrificed with a robust conscience to the general welfare.

\end{quotex}
One cannot possess the virtues, while remaining obsessed with lusts and the ``satisfaction of life". And where there is lawlessness, someone must lay down the law.



\flrightit{Posted on 2011-07-21 by Cologero }

\begin{center}* * *\end{center}

\begin{footnotesize}\begin{sffamily}



\texttt{Perennial on 2011-07-22 at 04:27 said: }

I may be speaking vainly here, but Savanarola always struck me as a bit of a revolutionary figure himself. He certainly had heroic aspects about him, but is he really a counter-revolutionary? Established history does not seem to view him as such. A revionism, perhaps?


\hfill

\texttt{Cologero on 2011-07-22 at 12:58 said: }

Although we never specifically mentioned the term counter-revolutionary, Perennial, perhaps you are inferring that. After all, the Renaissance is one of the earliest stages of the on-going revolution. Hence, any opposition to that is part of the counter-revolution. Roeder seems to agree with that, since he regards the Renaissance as a rejection of the ignorant Middle Ages, while Savonarola as a revert to that more Traditional age. Evola regards Savonarola as the last attempt to recover an earlier antique consciousness in opposition to the falsity of the Renaissance. (This is something known only by readers of Gornahoor in the Anglosphere, since it is a passing reference in one of his untranslated books.)

So, it is really up to you to show how Evola and Roeder misunderstand Savonarola and why we should lump him in, instead, with the likes of Luther, Robespierre, and Lenin.


\hfill

\texttt{Perennial on 2011-07-22 at 13:58 said: }

I am going to cite wikipedia here, not because I think it is a fantastic source, but because it normally reflects the general understanding of a subject:

``Savonarola's religious actions have been compared to those of the later 17th and 18th century Jansenists, although theologically many differences exist. Savonarola did not produce a theological doctrine on salvation, and faithfully adhered to even minor theological definitions of the papal Magisterium. However, Savonarola's call to simplicity in church interior and his rigorous moral stances have been compared to those of Jansenists. Also the insistence on the immediate danger of Hell and the fewness of the elect can be considered to be a similarity.

After Savonarola's death, a secret Catholic group known as the Piagnoni sprang up in Florence to preserve his memory, organized into a sort of Catholic guild. Franciscan friars were prominent among the Piagnoni, and they briefly re-appeared in 1527 when they once again overthrew the Medici, but through intervention of the Holy Roman Empire it was brought to an end in 1530 at the Battle of Gavinana and the Medici were restored to power.

Savonarola left many admirers throughout Europe, in particular among religiously pious humanists who valued his deep spiritual convictions. Erasmus is said to have refused to become a Protestant partly because of the influence of Savonarola. At the same time, he is considered by Protestants to be a forerunner of the Reformation because of his criticisms of the papacy. Savonarola was said to be an inspiration to Michelangelo and Martin Luther." 

If Savonarola was really building himself on ancient ideals, why was he an inspiration to protestants, humanists, revolutionaries and opponents of the Empire? I am just curious as to how he could have been so misunderstood by those he inspires, considering Joseph de Maistre and Evola are certainly not a party to any misunderstanding in our days. Perhaps it was his style, rather then substance? Would this be a warning to counter-revolutionary elements about avoiding the wrong appearance?


\hfill

\texttt{Matt on 2011-07-22 at 15:43 said: }

Perrenial,

Yes, it could very well be (and probably was) due to style. Since he criticised the pope at the time, that was probably the basis of his inspiration for the early protestants, as I'm sure any criticism of the papacy, regardless of the source, caught the admiration of the leaders of the Reformation. Does some of Savonarola's actions prove as a warning to present and future counter-revolutionaries? Possibly. I suppose it depends if one is willing to take those risks he took.


\hfill

\texttt{Cologero on 2011-07-22 at 23:11 said: }

In all fairness, that does not imply Evola's blanket endorsement of Savonarola's religious beliefs. Rather that the latter exhibits a certain spiritual attitude, or we could say, he exemplifies a specific ``race of the spirit". The point is that such judgments have little to do with someone's particular affiliation.


\hfill

\texttt{Perennial on 2011-07-23 at 02:25 said: }

So what we are really saying is, burn with the same fire as Savonarola, we just do not have to burn the same way. Savonarola stood as few men dare to stand. It seems, however,that his figure was ambiguous enough to be taken by the forces of subversion. Looking at my examples of de Maistre and Evola, for example, no ambiguity certainly exists there. Perhaps this is what I mean by being a warning to we of Tradition, that if we are to stand, we should stand clearly, marking our point of departure. Savonarola does not seem to have made his distinctions sufficiently clear, at least not to the forces of order, specifically the Empire. He certainly was not taken as a man of order either by the Empire, or the Church. As Christ says in Apocalypse: ``I know thy works, that thou art neither cold, nor hot. I would thou wert cold, or hot." (Apocalypse 3:15)Savonarola might not have been such an unclear figure if he had more forthright about what he was about, aside from a moral standpoint. But Savonarola may not have had a wider perspective, which may also be another weak point for him. In any case, his example is worthy of further study.


\hfill

\texttt{Cologero on 2011-07-23 at 10:05 said: }

The Luther connection is absurd. Savonarola rightfully criticized a particular pope for his lifestyle, not the papacy itself — a distinction you seem to understand from the comments on the priest-king. Luther's sensuousness was opposed to S., and while L. is often an inspiration to the neo-pagans (e.g., Chamberlain), he is the beginning of the end of Tradition in the West.

The only interesting point in that article is the Franciscan connection. S.'s own order, the Dominicans, abandoned him while the Franciscans felt a spiritual kinship. Compare S. to Francis, and more remotely to Augustine, the next link in the chain.

However, we are trying to avoid theological debates which only engage the mind. The issue is to understand the spirit of S., how does he view the world, how does he act in it, and so on, apart from any particular beliefs. S. has that dualist spiritual attitude to the world of nature, the attitude that Evola associates with the best of the pagans. It takes an imaginative meditation to put oneself in that position, a type of meditation we have described more than once. That is how we try to understand S., not by his ``influences" or those he ``influenced".

S. had visions and the gift of prophecy, which is an indication of an archaic consciousness.

As for your warning. The revolution will use any tool to further its end. My experience is that it is impossible to avoid the wrong appearance; most men do not think and cannot follow a train of thought. In particular those who are subjected to propaganda can only see ideas in terms of their indoctrination. I refer you to what we have written on astrology and on possession by elementals here and on the meditations blog. The first trial is not an optional exercise.


\hfill

\texttt{Visionsofglory14 on 2011-07-24 at 16:05 said: }

Perennial does make the good point that without proper projection, a Traditional force can be swept up and used as a tool of the Revolution. Certain Neo-Christian forces today are marked by members who have visions of the same character that Savonarola has: fiery and apocalyptic. Their doctrine and attitude toward the modern world reflect this, yet they have been turned into the spearhead of anti-Traditional forces at work today. A post made a few weeks back by Protestants attempting to rectify Tradition with their belief reflects this: both Luther and Calvin earnestly believed their work to a regeneration of Augustinianism, even though the effect of their doctrine was the subversion of it.


\hfill

\texttt{Cologero on 2011-07-24 at 16:35 said: }

Interesting thought. So you think a truck going backwards is the same as a truck going forwards? Hard to imagine.

In actuality, you have no idea what Luther and Calvin ``earnestly" believed, and I suspect it is quite different from your suppositions. A man is responsible for his own consciousness and a false consciousness cannot be justified by any amount of enthusiasm, earnestness, or invincible ignorance. He is also responsible for all the consequences of his acts, even those unintended.

As an exercise, you could ask those neo-Christians if they accept any of the main principles of Tradition. For example, ask one what the ``Logos" is and see if you get a coherent response.


\hfill

\texttt{Perennial on 2011-07-25 at 01:44 said: }

Actually, Cologero, that would be a humerous excercise. Do not tempt me to try it.


\hfill

\texttt{Gabe Ruth on 2013-08-29 at 14:44 said: }

Have you read Chesterton's little piece on the Savonarola? It seems to me that you are not fair to GKC. He did not mean that there have never been any Christians, but that modern man has rejected the thing without trying it, often because he thinks he's outgrown it. 

``The great deliverers of men have, for the most part, saved them from calamities which we all recognise as evil, from calamities which are the ancient enemies of humanity. The great law-givers saved us from anarchy: the great physicians saved us from pestilence: the great reformers saved us from starvation. But there is a huge and bottomless evil compared with which all these are fleabites, the most desolating curse that can fall upon men or nations, and it has no name except we call it satisfaction. Savonarola did not save men from anarchy, but from order; not from pestilence, but from paralysis; not from starvation, but from luxury. Men like Savonarola are the witnesses to the tremendous psychological fact at the back of all our brains, but for which no name has ever been found, that ease is the worst enemy of happiness, and civilisation potentially the end of man." 

Regarding Savonarola's being of the left, consider what effect he had on the world. He built something, and it was true. That it was destroyed by evil does not make it less so, and to doubt the man because he was used in propaganda of the left is very naive. He built something, and most astounding of all it seems he did it by example, rather than by force. His rhetoric inflamed the hearts of a city to turn to virtue, and his actions backed up his rhetoric.


\hfill

\texttt{Pickman on 2013-08-30 at 22:51 said: }

Why bother living at all is the essence of this article on ``life". Repent, repent and duly repent until your reach an old age of bitterness and resentment of never living at all. 

Live is my solution, live and not care for the reactions. Live by the heart and die by the sword my brethren and you will know ecstasies of spiritual heights beyond heaven itself! Joy my friends is the way and that only. Long live death!


\hfill

\texttt{Cologero on 2013-08-31 at 14:20 said: }

Yes, Pickman, \textit{Guadeamus igitur, juvenes dum sumus}\footnote{\url{http://www.youtube.com/watch?v=WfGXYkfLJ4s}} is the right philosophy for the young. But, I assure you, there is no bitterness in old age, because \textit{E' un bel giorno per morire}\footnote{\url{http://www.youtube.com/watch?v=X-tn5N8cw1k}}.


\end{sffamily}\end{footnotesize}
